\section{Kết luận}
\textcolor{red}{MARKED}
Nhóm đã thiết kế hệ thống mạng cho bệnh viện gồm trụ sở chính và hai cơ sở phụ: mô hình ba lớp ở trụ sở chính, mô hình lõi thu gọn ở cơ sở ĐBP và BHTQ, VLAN chia theo khoa/phòng ban, kết nối Internet Dual DSL cân bằng tải và kênh truyền riêng giữa các cơ sở. Phép tính lưu lượng cho thấy nhu cầu ra Internet vào khoảng 550~Mbps xuống và 100~Mbps lên sau tăng trưởng 20\%, được đáp ứng bởi hai đường 500~Mbps. [Tóm tắt kết quả mô phỏng và kiểm thử, điền sau.]

\subsection{Tóm tắt thiết kế}
Hệ thống dã thiết kế mạng cho 3 cơ sở với \textcolor{red}{\_SỐ\_} VLAN, kế hoạch địa chỉ dùng IPv4 dùng dải 10.0.0.0/8 chia theo khoa/phòng ban kết hợp tòa nhà/cơ sở phụ, mỗi subnet được tính với dự phòng tăng 20\% số máy trạm. \textcolor{red}{\_TÓM tẮT\_}.

\subsection{Đánh giá theo tiêu chí đề bài}
\begin{table}[H]
    \centering
    \small
    \renewcommand{\arraystretch}{1.4}
    \begin{tabularx}{\textwidth}{
        @{}
        >{\raggedright\arraybackslash}p{4cm}
        >{\raggedright\arraybackslash}X
        @{}
    }
        \toprule
        \textbf{Tiêu chí} & \textbf{Đánh giá} \\
        \midrule
        \textbf{Reliability}
        & 2 edge router, 2 firewall, 2 core switch; HA cho cấp cứu và hồi sức tích cực; server DR đặt ở BHTQ \\

        \textbf{Ease of upgrade}
        & VLAN và subnet theo khoa, dải IP dự phòng 20\%, khối địa chỉ riêng cho từng toà/site \\

        \textbf{Diverse support software}
        & Hạ tầng server hỗ trợ HIS, LIS, RIS, FIS, PIS; có DMZ cho email, DNS, web \\

        \textbf{Safety}
        & Phòng IT, data center, phòng cáp trung tâm tách riêng khỏi toà A và toà B \\

        \textbf{Network security}
        & Phân vùng VLAN, DMZ, firewall; thiết bị khách hàng không truy cập được PC trong LAN \\
        \bottomrule
    \end{tabularx}
    \caption{Đánh giá thiết kế mạng theo các tiêu chí kỹ thuật}
\end{table}

\subsection{Hạn chế}
\begin{itemize}
    \item Khảo sát chỉ ở mức tượng trưng (không khảo sát tại chỗ).
    \item Số thiết bị mạng ở trụ sở chính (khoảng 21 thiết bị) vượt mức tham chiếu 12 của đề, nhóm chọn ưu tiên đáp ứng HA và phân subnet theo khoa/phòng ban.
    \item Các thiết bị y tế/IoT chỉ nằm ở mức dự kiến, chưa nằm trong IP plan.
    \item Mỗi site phụ chỉ có một đường leased line về main site, nên vẫn còn SPOF.
    \item Kết nối WAN hiện chỉ sử dụng kênh truyền riêng (leased line), chưa triển khai các công nghệ WAN hiện đại như MPLS hoặc SD-WAN do giới hạn thời gian thực hiện. Mỗi cơ sở phụ chỉ có một đường truyền về trụ sở chính nên tồn tại điểm lỗi đơn (single point of failure) ở kết nối WAN.
    \item Tốc độ đường truyền tính theo trường hợp xấu nhất và giả định số thiết bị khách; thực tế DSL có thể thấp hơn.
    \item \textcolor{red}{Gateway redundancy} (HSRP/VRRP) chưa được triển khai, camera và thiết bị IoT mới tính ở mức dự kiến.
    \item \textcolor{red}{Các hạn chế phát sinh từ mô phỏng}.
\end{itemize}

\subsection{Định hướng phát triển}
Trong tương lai, hệ thống có thể được nâng cấp kết nối WAN lên MPLS hoặc SD-WAN kèm đường truyền dự phòng cho các cơ sở phụ, triển khai gateway redundancy (HSRP/VRRP) và hoàn thiện kế hoạch địa chỉ cho camera cùng các thiết bị IoT, y tế.
